\section{The actual verification frontier}\label{c:verify:frontier}

The gravity program now has an explicit chain of finite and resolved constructions: linearized constraint propagation, nonlinear detail return, spectral constraint closure, a real Lorentzian parent, a local canonical reduction, and spherical strong-field models. Their hypotheses are part of the results. They do not merge into one unrestricted theorem merely because their formulas can be assembled in one volume.

The formal source frontier remains Block 49. Its newest result concerns the maximal free Dirac operator on the physical torus. The accompanying preparation record did not run Lean and did not add a gravity module. The gravity monograph presents ordinary mathematical derivations and selected computational evidence. The title of this companion reflects that distinction.

There is nevertheless a genuine formal interface. The checked source contains the pointwise-derivation obstruction, matrix corner derivations, curvature commutators, exact projector coarsening, and a complete-metric reconstruction theorem. Eight modules and twenty-one selected declaration locations are mapped below. These provide reusable mathematics. They do not yet instantiate an Einstein--matter constraint algebra, a Grassmann-valued BV charge, or a nonlinear gravity evolution tower.

A new parity audit also found a substantive error in the inherited Lorentzian BRST presentation. The archived dual-spinor formula puts an odd Lorentz ghost on the right with the commuting-parameter sign. For odd fermionic matter and a left odd differential, that sign fails nilpotency. Section~\ref{c:verify:parity} gives the exact counterexample and the corrected right-action sign. The old numerical covariance checks use commuting spinors and parameters, so they could not detect this error. Their retained pass labels are preserved as historical evidence rather than promoted into authentication of the printed master equation. A second review gate leaves the physical massive-Dirac convention matching open, as explained in Section~\ref{c:verify:boundary}.

\subsection{Four meanings of a successful check}

An ordinary proof establishes a mathematical implication under its stated assumptions. A source-identity check establishes that a cited file is the exact retained archive member. A bounded computational check establishes the assertion actually implemented for its inputs and arithmetic. A Lean kernel check establishes a formal declaration in a specified environment. The accompanying record supplies the first three kinds in selected scopes and cites inherited formal declarations. It supplies no new kernel receipt.

The durable crosswalk is \path{editorial/gravity/verification_source_map.json}. It records seven central claim groups, sixty-nine distinct archive members, eight formal modules and twenty-one declaration locations. Every member has a SHA-256 digest and its original ZIP path. The source map is a route back to the evidence, not a count of independent physical laws.

\clearpage
\section{The central claim ledger}\label{c:verify:ledger}

The technical supplement identifies C08--C14 as the nonlinear and strong-field chain. Its chapter numbers follow the compiled Volume I order. This companion uses the exact source filenames as the stable locators, because those filenames retain earlier numbering.

\begin{longtable}{@{}L{.07\textwidth}L{.38\textwidth}L{.45\textwidth}@{}}
\toprule
ID & Result family & Scope needed to read the result\\
\midrule
\endhead
C08 & Homological BFV transfer & Nilpotent fine differential, admissible contraction, invertible return operator and a compatible Hamiltonian or cyclic structure.\\
C09 & Resolved Ashtekar algebra & Band-limited fields and smearings, sufficient headroom \(16K<N\), detail bands retained before aliasing.\\
C10 & Lorentzian classical BV parent & Declared Einstein--Cartan--Maxwell--Dirac action and consistent graded gauge conventions. The dual-spinor sign needs correction; physical massive-Dirac convention matching remains open.\\
C11 & Boundary BFV and time slab & Regular nondegenerate boundary reduction. ADM ghost coordinates require sufficient headroom \(32K<N\). The time-slab relation has separate gauge-fixing and solution hypotheses.\\
C12 & Second-class reduction and time step & Oriented regular time gauge, positive triad determinant, an invertible endpoint Dirac matrix, and an exactly solved implicit step.\\
C13 & Charged spherical mass coordinates & Spherical electrovac on a monotone areal-radius patch. Strong-field reconstruction has additional square-root and derivative lower bounds.\\
C14 & Scalar flux and trapped sphere & Spherical massless scalar in the declared flat-slice gauge. The formation benchmark uses regular excision and bounded numerical controls.\\
\bottomrule
\end{longtable}

The inherited seven ledgers contain 121 rows labelled pass: respectively 9, 9, 30, 16, 26, 19 and 12. They also retain one conditional-theorem row and two open rows in the parent and boundary ledgers. These counts are a source inventory. Some rows explicitly record analytic arguments with a zero residual. Those rows were not numerical evaluations of a complete BV charge. A zero in such a ledger must be read together with its method and note fields.

The current edition adds a bounded reproduction of the underlying second-class verifier, nine strong-field identity and archive-consistency checks, and seven exact BRST-parity checks. These three scopes are distinct. The parity checks include a successfully detected failure of the archived sign, so ``checks pass'' does not mean every inherited mathematical assertion survived unchanged.

\clearpage
\section{How the nonlinear route changed}\label{c:verify:chronology}

The early linearized ADM construction uses commuting derivative operators and a quadratic action. On that model, divergence identities propagate the linearized Hamiltonian and momentum constraints. The argument is valuable precisely because its cancellation can be inspected directly. It does not prove that the same identities survive arbitrary nonlinear pointwise multiplication.

The nonlinear chapter exposes the obstruction. In a finite pointwise algebra, an exact scalar-linear Leibniz derivation must vanish. Ordinary finite differences avoid being zero by failing that Leibniz identity. Contracting the defect into a particular constraint bracket may produce a visible nonzero witness, but one such witness is narrower than the universal algebraic obstruction. The current monograph preserves both levels of argument.

The next construction changes the algebra. Matrix corner derivations obey an exact product rule, while homological transfer retains coarse-to-detail-to-coarse returns. The transfer theorem assumes a consistent fine nilpotent theory and contraction data. It does not construct a microscopic four-dimensional gravity theory from those assumptions.

The resolved Ashtekar construction supplies a more concrete classical input. Trigonometric-polynomial products are allowed to occupy larger support bands. Below Nyquist, multiplication, differentiation, trace and canonical variation agree with the continuum polynomial calculation. This produces exact resolved bracket identities without making the complete finite pointwise algebra into a differential algebra.

The later Lorentzian chapters use a real first-order parent rather than imposing physical reality only after choosing one complex self-dual chart. The boundary construction then splits time and space. Its ledger still records second-class reduction as open at that stage. The following chapter provides a local regular time-gauge reduction. The earlier open row is therefore historical, while global reduction and degenerate strata remain open.

The discrete-gradient theorem supplies a finite-dimensional invariant-preserving step once a regular gauge patch and an exact step solution are available. The spherical electrovac and scalar chapters give additional reduced models and benchmarks. None of these later developments retrospectively proves a global Einstein--matter continuum limit or a quantum anomaly-cancellation theorem.



\clearpage
\section{What is already formal: the algebraic obstruction}\label{c:verify:formalalgebra}

The namespace \path{KomorebiPhysics.PointwiseNoGo} defines the pointwise algebra \(\iota\to\C\) and a complex-linear map \(D\) with the actual Leibniz property. Its theorem \path{finite_source_no_go} assumes a finite decidable index type and concludes \(D=0\). The stronger source theorem \path{derivation_eq_zero} uses coordinate idempotents without needing a finite-type hypothesis. Neither theorem concerns every possible finite-dimensional commutative algebra.

The proof is short enough to explain. A coordinate idempotent \(e_a\) satisfies \(e_a^2=e_a\). Leibniz gives
\[
 D e_a=2e_aD e_a.
\]
At coordinate \(a\), this says \(D e_a(a)=2D e_a(a)\); elsewhere it says \(D e_a(b)=0\). Thus \(D e_a=0\). Applying the rule to \(e_af=f(a)e_a\) forces \(D f(a)=0\) for every \(a\). This is the finite pointwise obstruction used by the gravity discussion.

The namespace \path{KomorebiPhysics.CornerDerivation} changes the product space to finite complex matrices and defines
\[
 \delta_KX=i(KX-XK).
\]
The theorem \path{delta_mul} proves
\[
 \delta_K(XY)=(\delta_KX)Y+X(\delta_KY).
\]
For Hermitian \(K\), \path{delta_preserves_hermitian} preserves Hermitian inputs. The theorem \path{trace_integration_by_parts} proves
\[
 \operatorname{tr}((\delta_KX)Y)
 =-\operatorname{tr}(X\delta_KY).
\]
These are actual algebraic identities for the source operators, rather than desired properties inserted into a structure.

In \path{KomorebiPhysics.CornerCurvature}, the curvature matrix is
\[
 F_{KL}=i[K,L].
\]
The declaration \path{derivation_commutator} proves
\[
 \delta_K\delta_LX-\delta_L\delta_KX=\delta_{F_{KL}}X.
\]
The commuting-generator hypothesis in \path{derivations_commute_of_generators_commute} is explicit. To use commuting corner derivatives in the linearized ADM proof, one must supply that hypothesis for the chosen generators. The existence of the theorem does not identify arbitrary matrices with physical spatial derivatives.



\clearpage
\section{What is already formal: balance and refinement}\label{c:verify:formalrefinement}

The formal refinement results act on projector atlases, finite sums and currents. In \path{KomorebiPhysics.Atlas}, \path{coarsen_energy} adds packet norm squares over a fiber, and \path{coarsen_current} adds the corresponding fine current entries. These results come from \path{KomorebiPhysics/Refinement.lean}. They ensure that exact coarse aggregation preserves the established balance law.

The namespace \path{KomorebiPhysics.RefinementTower} proves composition of aggregation. Its declarations \path{push_comp}, \path{incidence_comp} and \path{aggregateCurrent_comp} show that direct and iterated sums or incidence-matrix transfers agree. The current transfer has the form \(F\mapsto RFR^T\). This is exact path independence of that aggregation. It is not automatically path independence of nonlinear canonical gravity reductions.

The namespace \path{KomorebiPhysics.UniversalBalance} proves cancellation of internal antisymmetric fluxes. The theorem \path{source_boundary_sum} separates sources and boundary current; \path{hasDerivAt_total} removes internal transfer from the total derivative. To use this for spherical gravity, a separate derivation must identify the actual mass density, flux, differentiability and boundary conditions. The exact scalar mass-flux theorem remains that additional ordinary proof.

The strongest abstract continuum bridge here is in \path{KomorebiPhysics.LocalDiscreteGlobalContinuum}. An exact tower includes embeddings \(I_n\), level evolutions \(U_n\), and the assumption
\[
 U_{n+1}I_n=I_nU_n.
\]
The declaration \path{ExactRefinementTower.evolve_iterEmbed} propagates this relation through successive levels. A complete metric reconstruction then uses adjacent bounds
\[
 d(f_n,f_{n+1})\le b_n,
 \qquad \sum_{n=0}^{\infty}b_n<\infty.
\]
The declarations \path{SummableDefect.exists_limit} and \path{dist_limit_le_tail} give convergence and the tail estimate
\[
 d(f_n,f_\infty)\le\sum_{j=0}^{\infty}b_{n+j}.
\]
The theorem \path{ExactRefinementTower.ContinuumRealization.globalEvolution_error_bound} applies the same mechanism after levelwise evolution.

Gravity still needs its own state spaces, reconstruction maps, regular chart control, and summable evolution defects. A few-grid convergence fit does not supply that summable tower, and the abstract theorem does not construct it. This is a precise formal gap rather than an absence of useful formal progress.

\clearpage
\section{The spectral headroom gate}\label{c:verify:headroom}

For the resolved vacuum algebra, a field in \(\mathcal T_K\) has coordinatewise Fourier support at most \(K\). Products satisfy
\[
 \mathcal T_K\mathcal T_L\subseteq\mathcal T_{K+L}.
\]
If a product has support below \(N/2\), its sampled derivative agrees with the derivative of the trigonometric polynomial. The normalized grid trace is its exact zero Fourier mode. The rule is checked for the product actually used; a low support bound on each factor alone is insufficient.

The inherited C09 theorem counts the smearing among the band-\(K\) factors. Its largest canonical gradients have four factors, and a bracket has at most eight. The sufficient condition is therefore
\[
 8K<N/2,\qquad\text{or equivalently }16K<N.
\]
On this domain, the Hamiltonian bracket closes into the vector constraint, which is a diffeomorphism constraint plus a Gauss term. The densitized lapse has weight \(-1\); omitting that weight changes the mixed bracket.

The retained C09 computation uses an \(18^3\) grid with band 1 inputs, so its declared sufficient headroom holds. It records seven bracket residuals of order \(10^{-17}\), one support check and one aliasing negative control. The band 3 control violates the hypothesis and gives an absolute Hamiltonian-bracket defect \(3.288599\times10^{-8}\). This separates an observed resolved case from an observed aliased case.

A nonlinear Hamiltonian vector field can create new support. A fixed \(\mathcal T_K\) is not generally invariant under evolution. The construction therefore uses a filtered tower and refines before products leave the allowed headroom. Exact bracket evaluation below Nyquist is distinct from existence of a global, invariant finite-band flow.

The parent and boundary audit grids need additional care. Their retained ordinary covariance checks use \(N=12,K=1\) and a band 1 \(10^4\) grid respectively. Those sizes do not satisfy the sufficient whole-parent \(16K<N\) or whole-ADM-ghost \(32K<N\) bounds. Individual lower-degree identities can have smaller support requirements. Those computations cannot be cited as complete numerical verification of the full master charge under the stated sufficient bounds.



\clearpage
\section{The Lorentzian BRST parity correction}\label{c:verify:parity}

Fix a left odd derivation \(s\), so for homogeneous quantities
\[
 s(ab)=(sa)b+(-1)^{|a|}a(sb).
\]
Consider only the internal Lorentz spin sector. Let the ghost matrix \(\chi\) and the row spinor \(\bar\psi\) both be odd, with
\[
 s\chi=\chi^2.
\]
The archived formula places the ghost on the right and uses
\[
 s\bar\psi=-\bar\psi\chi.
\]
Its second derivative is
\[
 s^2\bar\psi
 =-\big((s\bar\psi)\chi-\bar\psi(s\chi)\big)
 =2\bar\psi\chi^2.
\]
That expression need not vanish. The sign appropriate for a commuting infinitesimal Lorentz parameter cannot be copied unchanged after moving an odd ghost through odd matter.

An exact finite witness uses Lorentz spin generators \(T_1=\gamma_{23}/2\), \(T_2=\gamma_{31}/2\) for the mostly-plus Clifford convention. Their commutator is \(-\gamma_{12}/2\ne0\). For independent odd generators \(\theta_1,\theta_2,\eta\), set
\[
 \chi=\theta_1T_1+\theta_2T_2,
 \qquad \bar\psi=\eta r,
 \qquad r\gamma_{12}\ne0.
\]
Then \(\chi^2=\theta_1\theta_2[T_1,T_2]\), and the archived sign gives a nonzero degree-three coefficient.

The consistent right-row convention is
\[
 s\bar\psi=+\bar\psi\chi.
\]
It gives \(s^2\bar\psi=\bar\psi\chi^2-\bar\psi\chi^2=0\). Equivalently, one may keep the negative dual-representation sign with all odd ghosts placed on the left in component formulas. Under this convention the mass bilinear also obeys
\[
 s(\bar\psi\psi)=(s\bar\psi)\psi-\bar\psi(s\psi)=0.
\]

The new exact check uses \(\mathbb Q(i)\), seven Grassmann generators and actual \(4\times4\) Clifford matrices. It verifies all 16,384 exterior-basis products and records seven passing checks, including the expected nonzero archived residual and the zero corrected residual. Its schema 2 manifest records a measured final run, actual input hashes before and after, and declared result hashes. This is a decisive finite counterexample to the printed pure-Lorentz nilpotency claim. It is not an independent verification of the complete diffeomorphism, electric, BV or boundary charge.



\clearpage
\section{Parent and boundary claims after the correction}\label{c:verify:boundary}

The real parent brings the tetrad, Lorentz connection, Maxwell auxiliary field and Dirac matter into one declared first-order action. Its ordinary commuting-parameter symmetry calculations remain useful evidence about curvature, torsion, covariance and the action density. The parity finding limits what follows from those checks: the printed graded differential needs correction before its nilpotency can support a classical master action.

There is a separate physical convention gate. The inherited signature is \(\eta=(-,+,+,+)\) and the Clifford relation is \(\{\gamma^I,\gamma^J\}=2\eta^{IJ}\). The flat operators \(i\gamma^I\partial_I\pm m\) have product \(-\Box-m^2\), so either real-mass sign gives \((\Box+m^2)\psi=0\). Plane waves therefore obey \(\omega^2=|k|^2-m^2\), rather than the conventional positive massive shell. The covariance checks do not test that dispersion. The displayed gauge algebra can be studied as a declared model, but matching its factors of \(i\), adjoint and mass sign to a physical massive Dirac action remains open. The free torus result from Block 49 does not repair this separate parent convention.

The classical BV construction has a clear dependency order. First establish a graded differential with \(s^2=0\). Then establish \(sS_0=0\) using the same parity and left/right conventions. Then insert the corresponding antifield terms and verify the chosen antibracket signs. An off-shell closed gauge algebra permits the usual minimal construction, but the actual component formulas must realize it consistently. The current monograph makes the dual-spinor repair explicit. The finite parity check certifies the local repair; it does not silently replace a full component audit.

The boundary construction has further assumptions. Its \(3+1\) split identifies multiplier constraints and candidate moment maps. The full normal moment-map statement remains conditional on embedding data, a projectable regular quotient and equivariance. Obtaining a canonical boundary bracket requires a regular nondegenerate reduction. ADM ghost coordinates are a field-dependent change of ghost variables, so the canonical pullback must include the conjugate momenta and their induced terms. Sufficient support headroom belongs to this complete coordinate construction, not merely to one multiplication test.

The adaptive time-slab relation is a relation defined by the resolved variational problem. Its charge-intertwining statement depends on regular gauge fixing and an exact or suitably controlled solution of that problem. It is not automatically an explicit, single-valued update map. The following local discrete-gradient construction is a separate finite-step result.

The later second-class chapter pulls the charge back only if the BFV vector field is tangent to the regular second-class surface. That tangency condition is an actual theorem hypothesis. It is not removed by presenting a solution map for the algebraic constraints.

Finally, the quantum master equation has an additional measure or BV-Laplacian term. A cyclic classical transfer does not automatically supply an anomaly-free regulated gravity--matter measure. The inherited C10 ledger correctly leaves the quantum master equation open. The present parity correction concerns the classical starting point and gives no quantum anomaly-cancellation result.



\clearpage
\section{The regular canonical reduction}\label{c:verify:reduction}

The local reduction works in an oriented time-gauge sector with
\[
 E^0=0,\qquad e=\det(e_a{}^i)>0.
\]
The gravitational simplicity pullback removes the spatial rotational-connection term from the symplectic potential and leaves the boost connection \(K_a{}^i\) paired with the densitized triad. The full antisymmetric internal-index sum is essential to the factor \(1/\kappa\).

Maxwell auxiliary elimination is algebraic on the nondegenerate sector. The electric momentum pairs with the spatial potential. The fermionic primary matrix is invertible for \(e\ne0\). Defining the half-density spinor \(\Psi=e^{1/2}\psi\) cancels the determinant-variation terms and gives the local reduced one-form
\[
\begin{aligned}
 \Theta_{\mathrm{red}}=\int_\Sigma\bigg[
 &\kappa^{-1}\widetilde E^a{}_i\,\delta K_a{}^i
 +\mathcal E^a\delta A_a\\
 &+\frac{i}{2}
   (\Psi^\dagger\delta\Psi-\delta\Psi^\dagger\Psi)
 \bigg]d^3x.
\end{aligned}
\]
The inherited convention gives
\[
 \{K_a{}^i(x),\widetilde E^b{}_j(y)\}_{D}
 =\kappa\delta_a{}^b\delta^i{}_j\delta^3(x-y),
 \qquad
 \{\Psi_A(x),\Psi_B^\dagger(y)\}_{D}
 =-i\delta_{AB}\delta^3(x-y).
\]

The torsion equation fixes the rotational connection algebraically and produces the axial-current interaction. In the parent signature and orientation conventions, the source gives a linear contorsion contribution \(-3\kappa J_5^2/8\) and the total on-shell action correction \(-3\kappa J_5^2/16\). Changing signature, orientation, gamma conventions or left/right fermionic derivatives requires rechecking those signs together.

The reduced chart still contains first-class rotation, electric, diffeomorphism and Hamiltonian constraints. ``Minimal'' means that redundant algebraic and fermionic-primary variables have been eliminated. It does not mean that all gauge orbits have been globally quotiented.

Degenerate triads are outside this regular chart. The inherited negative control sets \(e=0\) and obtains a rank-zero fermionic primary matrix. The local result therefore identifies a real boundary of applicability. It does not provide a stratified extension across that boundary or a global time-gauge section.

\clearpage
\section{Exact time steps and approximate solves}\label{c:verify:step}

At a finite resolved level, let \(J^T=-J\) be the constant canonical matrix, let \(H\) be an autonomous Hamiltonian, and let \(\Phi\) collect the constraints and regular gauge conditions. For endpoints \(z_0,z_1\), the average-vector-field discrete gradient satisfies
\[
 F(z_1)-F(z_0)
 =\overline\nabla F(z_0,z_1)^T(z_1-z_0).
\]
Let \(G_\Delta\) have the discrete gradients of the components of \(\Phi\) as its rows. Assume the endpoint matrix
\[
 \Delta_\Delta=G_\Delta JG_\Delta^T
\]
is invertible. Define
\[
 J_D^\Delta
 =J-JG_\Delta^T\Delta_\Delta^{-1}G_\Delta J.
\]
Because \(\Delta_\Delta\) and its inverse are skew, \(J_D^\Delta\) is skew. Direct multiplication gives \(G_\Delta J_D^\Delta=0\). Thus every exact solution of
\[
 z_1-z_0=\Delta t\,J_D^\Delta\overline\nabla H
\]
preserves \(\Phi\) and \(H\). This implication is finite-dimensional linear algebra plus the exact discrete-gradient identity.

The approximate-solve distinction can also be quantified. If the implemented endpoint has residual
\[
 r=z_1-z_0-\Delta t\,J_D^\Delta\overline\nabla H,
\]
while the discrete gradients and inverse are exact, then
\[
 \Phi(z_1)-\Phi(z_0)=G_\Delta r,
 \qquad
 H(z_1)-H(z_0)=\overline\nabla H^Tr.
\]
Consequently,
\[
 \|\Delta\Phi\|\le\|G_\Delta\|\|r\|,
 \qquad
 |\Delta H|\le\|\overline\nabla H\|\|r\|.
\]
Quadrature and floating-point errors add their own terms. Small solver residuals are evidence for small invariant drift in a controlled patch; they are not exact preservation in finite arithmetic.

Local step existence uses a regular neighborhood and sufficiently small \(\Delta t\). Long-time bounds require a coercive Hamiltonian and continued membership in the patch. An infinite tower additionally needs level-independent inverse-matrix, derivative and coercivity bounds. Those uniform gravity--matter estimates remain open.

\clearpage
\section{The new bounded second-class reproduction}\label{c:verify:run}

The accompanying record reran the underlying source program \path{Resolved_Second_Class_Reduction_Verification.py} in an isolated evidence directory, with unchanged source bytes. The seed is 20260811. The program combines symbolic torsion identities with floating-point pullbacks, Clifford identities, rank checks and one constrained-particle trajectory.

The trajectory is a six-dimensional surrogate with
\[
 q\cdot q=1,\qquad q\cdot p=0.
\]
It uses 1,000 steps of size 0.1. It is not a simulated gravity--matter trajectory. The final run returned 26 passes, 4 information entries and 0 failures. The retained and new values agree. Representative values are

\begin{center}
\begin{tabular}{@{}lr@{}}
\toprule
Bounded assertion & New value\\
\midrule
Energy drift & \(7.771561\times10^{-16}\)\\
Constraint drift & \(3.885781\times10^{-15}\)\\
Maximum implicit-solve residual & \(1.092187\times10^{-16}\)\\
RK4 energy-drift negative control & \(1.396105\times10^{-6}\)\\
RK4 constraint-drift negative control & \(2.594613\times10^{-6}\)\\
Gravity symplectic-pullback residual & \(1.110223\times10^{-16}\)\\
Half-density identity residual & \(3.662053\times10^{-15}\)\\
Exact on-shell contorsion coefficient & \(-3/16\)\\
\bottomrule
\end{tabular}
\end{center}

The implementation thresholds are \(10^{-10}\) for selected residuals and \(10^{-8}\) for the negative-control drift. The exact symbolic coefficient is compared as an exact string after the symbolic calculation. The rank tests and determinant tests have their own conditions. These heterogeneous checks should not be compressed into a single universal residual.

The manifest records Python 3.12.14, NumPy 2.3.5, SciPy 1.17.0 and SymPy 1.14.0, the command, actual runtime, source hash, output hashes and the finite bounds. It validates as a legacy version 1 computation record. Version 1 validation does not by itself enforce machine-level resource caps or a complete dependency graph. This run adds its input hash and confirms it remained unchanged, but the interpretation stays ``result reproduced; implementation not independently validated.''



\clearpage
\section{The electrovac verification boundary}\label{c:verify:electrovac}

The charged spherical model defines the mass coordinate
\[
 \mathcal M=\frac R2\left[
 1-\left(\frac{R'}\Lambda\right)^2
 +\left(\frac{P_\Lambda}R\right)^2\right]
 +\frac{Q^2}{2R}.
\]
Its exact identity expresses \(\mathcal M'\) as a combination of the Hamiltonian, radial momentum and electric constraints. Recovering those constraints from the triangular mass-coordinate system needs nonvanishing coefficients on a monotone regular patch. The theorem is not an algebraic equivalence at every degenerate configuration.

The finite-range benchmark evolves a radius variable and reconstructs the other canonical fields from mass and charge. That projection can make reduced mass-shell residuals close to roundoff. Standard ADM residuals evaluated with independent finite-difference derivatives remain nonzero and converge. A roundoff-level projected residual must therefore not be advertised as exact agreement of all discrete ADM operators.

The retained benchmark uses mass 1, charge 0.8, a radial interval from 0.6 to 6, a smooth coordinate deformation, final time 0.4 and CFL factor 0.10. Its refinement grids have 64, 128, 256 and 512 nodes. The inherited ledger has 19 passes, including symbolic identities, fourth-order reconstruction fits, horizon-coordinate consistency, a stability diagnostic and an unprojected negative control.

The radius evolution stencil is finite range. A radial constraint prefix solve and the reconstructed geometry have a different locality property. The full calculation is therefore not a proof that every field update in the reduced gravity solver is strictly finite range.

The new bounded strong-field audit calls the source's symbolic electrovac identities and obtains exact zero. It also checks an independently written derivative identity needed by the regular source estimate. It reads the existing convergence and check data without rerunning the full evolution. The distinction is recorded in its manifest.

The uniform stability route still needs lower bounds for the radius derivative and the square-root radicand, controlled coefficients, boundary data and a suitable common norm. An observed small spread of amplification over the tested grids is useful evidence for that route. It is not a theorem uniform over all refinement levels or over nonspherical data.



\clearpage
\section{The scalar-collapse evidence}\label{c:verify:scalar}

The massless spherical scalar uses the flat-slice Painlev\'e--Gullstrand chart with \(v=\sqrt{2m/r}\). Its exact reduced identities include
\[
 \partial_rm=r^2\left[\frac12(\Pi^2+\Phi^2)+v\Pi\Phi\right],
\]
\[
 \partial_tm=r^2\sigma(\Pi+v\Phi)(\Phi+v\Pi).
\]
The symbolic compatibility calculation relates the time derivative of the radial mass constraint to the radial derivative of this flux. It uses the actual lapse constraint and scalar evolution on the regular region \(r>0\), \(\sigma>0\), \(v>0\). An extension to vacuum points \(v=0\) is separate. The continuum proof does not establish consistency of thresholded or clipped numerical reconstruction there. The outgoing null expansion vanishes when \(2m/r=1\) in the declared chart. This identifies a marginally trapped sphere, not a global event horizon.

The retained strong-packet table has seven grids from 600 to 2,400 nodes. Every listed strong run reaches the trapped-sphere criterion. A weak packet stays untrapped through time 5. The event-coordinate fits use the six finest grids, from 800 to 2,400, and give orders 2.173853 for time and 2.084395 for radius. The600-to800 event time is nonmonotone. The fitted evidence is therefore not an all-grid monotonicity or asymptotic-order theorem.

At the finest grid, the interpolated event values are
\[
 t_H=4.6918209458,\qquad r_H=3.1134115480,
 \qquad m_H=r_H/2.
\]
The stored field diagnostics come from the later stopping time 4.6920696969. Its gap from the interpolated event is approximately \(2.48751\times10^{-4}\). The lapse, curvature, outer mass and event-field profile should be identified with that stored field time unless explicitly interpolated. Combining all of them under one exact event-time label loses this distinction.

The finest reported relative residuals are approximately \(2.2732\times10^{-6}\) for the radial mass constraint, \(2.8077\times10^{-6}\) for the lapse constraint, and \(1.2697\times10^{-3}\) for the independent finite-time mass flux. The outer mass change is about 0.2624 percent. These are controlled numerical residuals, not exact conservation in the discretized evolution.

The spatial characteristic stencil has radius two. Geometry is reconstructed by triangular radial prefix solves at every RK stage. Those solves are not strictly finite-range maps. The benchmark has excision rather than regular-center reflection, and it does not establish a critical exponent, singularity continuation, nonspherical stability or a complete \(3+1\) solver. The new symbolic audit reproduced the reduced identities but performed zero full collapse-evolution steps.

\clearpage
\section{Reproducing the evidence without confusing the wrappers}\label{c:verify:reproduce}

The edition ZIP preserves the exact earlier archives in \path{recovered_sources/}. To inspect an inherited gravity claim, extract the Volume I source archive into a new working directory. Its root is \path{Projector_Refinement_Physics_Volume_I_Source/}. Use that extracted root as the working directory for its source commands. This avoids writing regenerated outputs over the preserved archive or over the edition evidence.

The technical supplement's labels V1E47 and V1E54--59 are code numbers, not chapter numbers. V1E54--56 reuse retained JSON by default. They invoke the underlying verifier only if its JSON is missing or the respective \path{RERUN_ASHTEKAR}, \path{RERUN_LORENTZIAN_BV} or \path{RERUN_HEAVY_BFV} flag is 1. V1E57--59 convert existing results into ledger rows. Re-executing a converter is not a new evolution or symbolic reproduction.

The inherited archive does not pin a Python dependency environment. Ashtekar verification uses JAX with its 64-bit option enabled; reduction uses NumPy, SciPy and SymPy; scalar evolution optionally uses Numba. Historical library versions are not recorded in the selected outputs. The current reproduction records its own versions and hashes.

The distinction is especially visible in V1E57. It reads \path{resolved_second_class_reduction_checks.csv} and selects existing pass rows. Several standardized residual fields are set to zero while the actual recorded value is retained in the note. For example, the RK4 drift is a successfully detected negative control, not a zero-drift numerical result. The raw source-key value is the meaningful quantity.

To bypass the three retained-JSON caches, run the underlying programs from the extracted root. The nonlinear transfer code and second-class verifier can also be run directly:
\begin{lstlisting}
python code/47_nonlinear_bfv_bv_polyphase_bridge.py
python verification/Resolved_Ashtekar_Closure_Verification.py
python verification/Resolved_Lorentzian_BV_Parent_Verification.py
python verification/Resolved_Lorentzian_3plus1_BFV_Verification.py
python verification/Resolved_Second_Class_Reduction_Verification.py
\end{lstlisting}
The second-class program writes its JSON next to itself; the retained publication copy is in \path{data/}. Compare payloads rather than assuming identical locations. The accompanying record copies the unchanged verifier into a dedicated evidence directory and records its output there.

The underlying strong-field programs accept an explicit output directory:
\begin{lstlisting}
python verification/Finite_Range_Spherical_Electrovac_Verification.py \
  --output /absolute/path/to/electrovac_run
python verification/Finite_Range_Spherical_Scalar_Collapse_Verification.py \
  --output /absolute/path/to/scalar_run
\end{lstlisting}
The scalar verifier also exposes run controls through its command-line parser. Changing grids, time bounds, packet amplitude or controls changes the tested assertion. Record them rather than comparing an altered run to the archived seven-grid result as if it were identical.

For the current edition's narrower checks, use its root directory:
\begin{lstlisting}
python editorial/gravity/collect_verification_sources.py
python editorial/gravity/run_second_class_verification.py
python tools/verify_strong_field_identities.py
python tools/verify_gravity_brst_parity.py
\end{lstlisting}
The first command only reads ZIP members and regenerates the crosswalk. The other commands are bounded and record their own distinct claims. None invokes Lean.

\clearpage
\section{The inherited Lean environment and its limits}\label{c:verify:environment}

The formal archive is \path{Komorebi_Physics_Lean_v0.49.0_Checked_Block_49.zip}. Its toolchain is \path{leanprover/lean4:v4.34.0-rc2}, and the Mathlib revision in \path{lakefile.toml} is
\begin{center}
\path{aeb72f1e9981b6d35b426b16dca7759be3e7ba06}.
\end{center}
The retained package metadata still says version 0.48.0. That is a metadata lag in the Block49 archive, not evidence of a different theorem or a new compilation in this account. Both files are preserved and hashed.

The companion's source search covers 146 \path{KomorebiPhysics/*.lean} members in that archive. It searches source text for the direct gravity terms and examines the eight selected modules. It identifies no direct C08--C14 gravity declaration in that scope. This is a scoped source finding. It is not a semantic proof about every imported dependency, and it does not use the older global 142-module inventory as a current whole-project count.

The newest selected formal source, \path{Relativistic3DWeightedDiracDomain.lean}, is included to fix the actual frontier. In namespace \path{KomorebiPhysics.Relativistic3DWeightedDomain}, \path{mem_physicalDiracDomain_iff_summable_h1} identifies the maximal free physical domain with the first-order weighted Fourier condition. The declarations \path{physicalDiracPMap_selfAdjoint} and \path{physicalDiracPMap_isClosed} establish the self-adjoint and closed free operator. These are free torus statements. They do not include a dynamical tetrad, torsion, lapse, shift or gravity constraint surface.

The previous gauge companion already records retained-receipt exceptions: the Block 27 cumulative-build log contains a Git exit-code 128 failure inconsistent with the receipt, and six reported Blocks 36 and 40 logs are missing. This gravity edition preserves that finding and does not convert it into a mathematical counterexample. The direct source crosswalk here is independent of claiming that every inherited cumulative build receipt is complete.

A fresh formal reproduction would require the exact Lean and Mathlib environment and the archived release-specific commands. It must record successful source compilation and an axiom/trust audit. Merely running XeLaTeX, searching for a theorem name or observing a zero Python residual does not provide a new Lean check. No such execution was performed here.



\clearpage
\section{What remains to prove}\label{c:verify:remaining}

The closest new formal target is the finite discrete-gradient algebra. It would define exact discrete gradients, prove the endpoint tensor is skew and annihilates the constraint rows, and derive exact invariant preservation for a solved step. It should keep existence of the implicit solution separate. This target is narrow enough that its assumptions and output can be named without claiming a full gravity solver.

The graded parent needs a complete parity-consistent component audit. The newly detected Lorentz dual-spinor defect is repaired locally. The next check must cover the same ghost-placement convention in the diffeomorphism and electric sectors, the parent action, antifield derivatives, boundary moment maps and ghost momenta. A pure-Lorentz witness does not discharge those dependencies.

The matter realization also needs physical massive-Dirac convention matching. The inherited mostly-plus Clifford matrices with flat operators \(i\gamma^I\partial_I\pm m\) give the wrong sign for the conventional massive shell. A repair must treat the kinetic factor, adjoint, mass convention and any induced torsion signs together; covariance alone cannot select it.

The resolved vacuum algebra needs an explicit formal Fourier-support calculus, canonical bracket, constraint densities, and a verified translation of the continuum self-dual identities. The support theorem must apply to all intermediate products and generated smearings. Uniform propagation of headroom is a separate dynamics problem.

The Lorentzian canonical chart needs a formal treatment of nondegeneracy, graded Dirac brackets, torsion elimination and boundary-charge tangency. Extending the chart across degenerate triads or constructing a global gauge quotient requires additional geometry. The local time-gauge formulas do not supply that extension.

The infinite tower needs level-independent regularity, coercivity and inverse-Dirac-matrix bounds; compatible evolution and reconstruction maps; and explicit summable defects in a common complete space. Only then can the already formal metric-limit theorem be instantiated for an Einstein--matter trajectory.

The spherical numerical route needs independent solver reproduction, a proof of stability and consistency for the complete projected system with its boundaries, and a quantified event-location limit. A global trapped-surface theorem would require additional continuum existence and geometric estimates. Regular-center evolution, critical collapse, singularity formation, event horizons and nonspherical stability remain further targets.

The quantum route needs a regulator, measure and quantum master equation with anomaly control. It is not completed by the classical transfer formula or the exact finite nonlinear example. Empirical validation and a comparison to independent physical data are further tasks beyond this mathematical verification companion.

\clearpage
\section{A reader's verification contract}\label{c:verify:contract}

For any result in this gravity set, first identify the variables and the algebra. A Hermitian matrix derivation, a filtered Fourier field, a real tetrad, a half-density spinor and a spherical mass coordinate are different objects. Transfer between them requires a declared map and preserved structure.

Next read the domain. The essential gates include Fourier headroom, positive triad determinant, regular gauge fixing, endpoint-matrix invertibility, monotone radius, positive reconstruction radicand and continued membership in an excised chart. Some are pointwise conditions, some are neighborhood conditions, and some must be uniform over levels. The proof status changes when one of these conditions is removed.

Then identify the evidence type. For a symbolic zero, inspect the expression and arithmetic. For a floating residual, inspect the normalization, threshold, grid, seed and negative control. For an inherited pass row, inspect its method and whether the executable actually computed the claim. For a formal declaration, inspect the types and hypotheses rather than the physical phrase in its filename.

Finally trace dependencies. A classical master claim needs a consistent graded differential. A reduced master claim needs regular reduction and tangency. Exact invariant preservation needs an exactly solved step. A continuum claim needs uniform estimates and compatible reconstruction. A quantum claim needs its measure and regulator. The new BRST counterexample shows why those dependencies must remain explicit even when a large inherited ledger contains many passes.

The accompanying record provides a coherent gravity monograph, a Python visual atlas and this companion as three views of the same source. The manuscript gives the ordinary arguments, the atlas shows selected geometry and bounded numerical behavior, and the companion makes the remaining formal bridge and the evidence limits inspectable.

\subsection{Portable evidence index}

\begin{itemize}
\item \path{editorial/gravity/verification_source_map.json}: original member paths, hashes, seven claim groups and selected Lean locations.
\item \path{evidence/gravity/second_class_computation_manifest.json}: fresh bounded underlying reduction reproduction.
\item \path{evidence/gravity/strong_fields_identity_checks/}: symbolic strong-field identities and retained-table consistency checks.
\item \path{evidence/gravity/brst_parity_results.json}: exact negative control and corrected pure-Lorentz dual action.
\item \path{editorial/gravity/verification_layout_review.json}: final companion layout and proofreading scope.
\end{itemize}
